% GENERATED FILE. Edit canonical lessons, then run npm run generate:books.
% canonical-source-hash: fnv1a64:fa3e984f80f333d3
% canonical-lessons: ZH-C59-chuan, ZH-C59-xie, ZH-C59-maozi, ZH-C59-kuzi, ZH-C59-piaoliang

\chapter{Clothes}
\label{ch:zh-clothes}
\hlchaptermodality{59}

\begin{chapteropening}
\textbf{By the end of this chapter:} \emph{I can say what I wear: shoes, a hat and trousers, and that they look nice.}

Complete the last lesson of chapter 59: I can say what I wear: shoes, a hat and trousers, and that they look nice.
\end{chapteropening}

\section[chuān]{chuān — to wear; to put on}

\label{lesson:ZH-C59-chuan}



\begin{quote}
Before the new one: say the Mandarin for cup, then the Mandarin for clothes.
\end{quote}

\subsection*{You'll want to know: chuān}
\emph{chuān} — \textquotedblleft{}to wear; to put on\textquotedblright{}.

\begin{quote}
\emph{míngtiān nǐ chuān shénme} — What are you wearing tomorrow?
\end{quote}

Its characters come later; for now it is a word for the ear.

\begin{cousinweb}
\emph{chuān} is inherited from Middle Chinese \emph{tsyhwen} (Guangyun, in Baxter's transcription), glossed 'to go through' and 'a hole'. Wearing is not given there; putting arms and legs through clothes is the usual explanation, but it is not settled.
\end{cousinweb}

\begin{grammarlens}[title={What you've built}]
One more word for everyday answers.
\end{grammarlens}

\subsection*{Your turn}
\begin{itemize}
\raggedright
  \item \emph{Say it:} \emph{chuān}
  \item \emph{Say it:} \emph{chuān}, once more
  \item \emph{Say it:} \emph{míngtiān nǐ chuān shénme} — What are you wearing tomorrow?
  \item \emph{Recall:} say the Mandarin for cup, then the Mandarin for clothes, then say \emph{chuān} again
\end{itemize}

\begin{tcolorbox}[breakable,colback=teal!4,colframe=teal!35!black,title={Before you move on}]
What does chuān mean? (\textquotedblleft{}To wear; to put on\textquotedblright{}.) Say it once more.
\end{tcolorbox}
\section[xié]{xié — shoe}

\label{lesson:ZH-C59-xie}



\begin{quote}
Before the new one: say the Mandarin for clothes, then the Mandarin for to wear.
\end{quote}

\subsection*{You'll want to know: xié}
\emph{xié} — \textquotedblleft{}shoe\textquotedblright{}.

\begin{quote}
\emph{wǒ de jiǎo hěn dà xié hěn xiǎo} — My feet are big and the shoes are small.
\end{quote}

Its characters come later; for now it is a word for the ear.

\begin{cousinweb}
\emph{xié} is inherited from Middle Chinese \emph{heaj} (Guangyun, in Baxter's transcription), glossed 'shoe'.
\end{cousinweb}

\begin{grammarlens}[title={What you've built}]
One more word for everyday answers.
\end{grammarlens}

\subsection*{Your turn}
\begin{itemize}
\raggedright
  \item \emph{Say it:} \emph{xié}
  \item \emph{Say it:} \emph{xié}, once more
  \item \emph{Say it:} \emph{wǒ de jiǎo hěn dà xié hěn xiǎo} — My feet are big and the shoes are small.
  \item \emph{Recall:} say the Mandarin for clothes, then the Mandarin for to wear, then say \emph{xié} again
\end{itemize}

\begin{tcolorbox}[breakable,colback=teal!4,colframe=teal!35!black,title={Before you move on}]
What does xié mean? (\textquotedblleft{}Shoe\textquotedblright{}.) Say it once more.
\end{tcolorbox}
\section[màozi]{màozi — hat; cap}

\label{lesson:ZH-C59-maozi}



\begin{quote}
Before the new one: say the Mandarin for to wear, then the Mandarin for shoe.
\end{quote}

\subsection*{You'll want to know: màozi}
\emph{màozi} — \textquotedblleft{}hat; cap\textquotedblright{}.

\begin{quote}
\emph{wǒ de tóu hěn dà màozi tài xiǎo le} — My head's big; the hat is too small.
\end{quote}

Its characters come later; for now it is a word for the ear.

\begin{cousinweb}
\emph{màozi} is built from mào and the noun ending zi, both inherited from Middle Chinese, \emph{mawH} and \emph{tsiX} (Guangyun, in Baxter's transcription); the dictionary glosses the first as a covering for the head.
\end{cousinweb}

\begin{grammarlens}[title={What you've built}]
One more word for everyday answers.
\end{grammarlens}

\subsection*{Your turn}
\begin{itemize}
\raggedright
  \item \emph{Say it:} \emph{màozi}
  \item \emph{Say it:} \emph{màozi}, once more
  \item \emph{Say it:} \emph{wǒ de tóu hěn dà màozi tài xiǎo le} — My head's big; the hat is too small.
  \item \emph{Recall:} say the Mandarin for to wear, then the Mandarin for shoe, then say \emph{màozi} again
\end{itemize}

\begin{tcolorbox}[breakable,colback=teal!4,colframe=teal!35!black,title={Before you move on}]
What does màozi mean? (\textquotedblleft{}Hat; cap\textquotedblright{}.) Say it once more.
\end{tcolorbox}
\section[kùzi]{kùzi — trousers}

\label{lesson:ZH-C59-kuzi}



\begin{quote}
Before the new one: say the Mandarin for shoe, then the Mandarin for hat.
\end{quote}

\subsection*{You'll want to know: kùzi}
\emph{kùzi} — \textquotedblleft{}trousers\textquotedblright{}.

\begin{quote}
\emph{zhè shì wǒ gēge de kùzi} — These are my brother's trousers.
\end{quote}

Its characters come later; for now it is a word for the ear.

\begin{cousinweb}
\emph{kùzi} is built from kù and the noun ending zi, both inherited from Middle Chinese, \emph{khuH} and \emph{tsiX} (Guangyun, in Baxter's transcription). The dictionary lists the first under a variant form of the character and explains it, through an older book, as a garment for the shins.
\end{cousinweb}

\begin{grammarlens}[title={What you've built}]
One more word for everyday answers.
\end{grammarlens}

\subsection*{Your turn}
\begin{itemize}
\raggedright
  \item \emph{Say it:} \emph{kùzi}
  \item \emph{Say it:} \emph{kùzi}, once more
  \item \emph{Say it:} \emph{zhè shì wǒ gēge de kùzi} — These are my brother's trousers.
  \item \emph{Recall:} say the Mandarin for shoe, then the Mandarin for hat, then say \emph{kùzi} again
\end{itemize}

\begin{tcolorbox}[breakable,colback=teal!4,colframe=teal!35!black,title={Before you move on}]
What does kùzi mean? (\textquotedblleft{}Trousers\textquotedblright{}.) Say it once more.
\end{tcolorbox}
\section[piàoliang]{piàoliang — pretty; good-looking}

\label{lesson:ZH-C59-piaoliang}



\begin{quote}
Before the new one: say the Mandarin for hat, then the Mandarin for trousers.
\end{quote}

\subsection*{You'll want to know: piàoliang}
\emph{piàoliang} — \textquotedblleft{}pretty; good-looking\textquotedblright{}.

\begin{quote}
\emph{tā de yǎnjing hěn piàoliang} — She has beautiful eyes.
\end{quote}

Its characters come later; for now it is a word for the ear.

\begin{cousinweb}
\emph{piàoliang} is built from piào and liang, both inherited from Middle Chinese, \emph{phjiewH} and \emph{ljangH} (Guangyun, in Baxter's transcription). The dictionary glosses the first as beating silk floss in water and the second as 'bright'; how the pair came to mean pretty is not settled.
\end{cousinweb}

\begin{grammarlens}[title={What you've built}]
One more word for everyday answers.
\end{grammarlens}

\subsection*{Your turn}
\begin{itemize}
\raggedright
  \item \emph{Say it:} \emph{piàoliang}
  \item \emph{Say it:} \emph{piàoliang}, once more
  \item \emph{Say it:} \emph{tā de yǎnjing hěn piàoliang} — She has beautiful eyes.
  \item \emph{Recall:} say the Mandarin for hat, then the Mandarin for trousers, then say \emph{piàoliang} again
\end{itemize}

\begin{tcolorbox}[breakable,colback=teal!4,colframe=teal!35!black,title={Before you move on}]
What does piàoliang mean? (\textquotedblleft{}Pretty; good-looking\textquotedblright{}.) Say it once more.
\end{tcolorbox}
